\documentclass[12pt,a4paper]{article}
\usepackage[margin=25mm, headheight=15pt, headsep=10pt]{geometry}
\usepackage{fontspec}
\usepackage[table]{xcolor} % Note: table option is required for rowcolors
\usepackage{titlesec}
\usepackage{tocloft}
\usepackage{fancyhdr}
\usepackage{booktabs}
\usepackage{tcolorbox}
\tcbuselibrary{skins, breakable}
\usepackage[colorlinks=true, linkcolor=emerald, urlcolor=emerald, bookmarks=true]{hyperref}

\definecolor{emerald}{RGB}{0,102,51}
\definecolor{darkred}{RGB}{139,0,0}
\definecolor{lightgray}{RGB}{245,245,245}

\usepackage[nil]{babel}
\babelprovide[import, main]{bengali}
\babelprovide[import]{arabic}
\babelfont{rm}[Renderer=HarfBuzz,Scale=1.0]{Noto Serif Bengali}
\babelfont{sf}[Renderer=HarfBuzz]{Noto Sans Bengali}
\babelfont{tt}[Renderer=HarfBuzz]{Noto Sans Bengali}
\babelfont[arabic]{rm}[Renderer=HarfBuzz,Scale=1.1]{Noto Naskh Arabic}

% Section styling
\titleformat{\section}{\Large\bfseries\color{emerald}}{\thesection.}{0.5em}{}
\titleformat{\subsection}{\large\bfseries\color{emerald!85!black}}{\thesubsection.}{0.5em}{}
\titleformat{\subsubsection}{\normalsize\bfseries\color{emerald!70!black}}{\thesubsubsection.}{0.5em}{}
\titlespacing*{\section}{0pt}{18pt}{8pt}

% Fancy Header/Footer setup
\pagestyle{fancy}
\fancyhf{} % clear all header and footer fields
\fancyhead[L]{\color{emerald}\small\textbf{\nouppercase{\leftmark}}}
\fancyfoot[L]{\scriptsize\color{black!50}{\lat{AI}}-সংকলিত, আলেম-যাচাইকৃত নয়}
\fancyfoot[C]{\color{emerald!70!black}\thepage}
\renewcommand{\headrulewidth}{0.4pt}
\renewcommand{\headrule}{\hbox to\headwidth{\color{emerald}\leaders\hrule height \headrulewidth\hfill}}
\renewcommand{\footrulewidth}{0pt}

% Custom boxes
% 1. Ayat/Hadith Box
\newtcolorbox{ayatbox}[1][]{
    enhanced, breakable, 
    colback=emerald!5, 
    colframe=emerald, 
    boxrule=0.5pt, 
    arc=3pt, 
    left=10pt, right=10pt, top=10pt, bottom=10pt,
    #1
}

% 2. Quote/Translation Box
\newtcolorbox{quotebox}[1][]{
    enhanced, breakable,
    frame hidden,
    colback=lightgray,
    borderline west={3pt}{0pt}{emerald},
    left=15pt, right=10pt, top=10pt, bottom=10pt,
    sharp corners,
    #1
}

% 3. AI-generated notice box
\newtcolorbox{warnbox}[1][]{
    enhanced, breakable,
    colback=red!4,
    colframe=darkred,
    boxrule=0.8pt,
    arc=3pt,
    left=10pt, right=10pt, top=10pt, bottom=10pt,
    #1
}

% Commands
\newcommand{\arab}[1]{\foreignlanguage{arabic}{#1}}
\newcommand{\kib}[1]{\textcolor{emerald}{\textbf{#1}}}

% Latin (English) fallback: Noto Serif Bengali has NO Latin glyphs
\newfontfamily\latfont[Renderer=HarfBuzz]{Noto Serif}
\newcommand{\lat}[1]{{\latfont #1}}

\setlength{\parskip}{5pt}
\setlength{\parindent}{0pt}

% Fix for missing bullet in Noto Serif Bengali
\renewcommand{\labelitemi}{$\bullet$}



\begin{document}
\begin{center}{\Large\bfseries\color{emerald} 09-উমার-শামে-সাধারণ-পোশাক}\end{center}
\vspace{1em}
\section*{ঘটনা ৯: উমার (রাঃ) — শামে সাদা পোশাকে ইসলামের সম্মান}

\begin{warnbox}
 \textbf{গুরুত্বপূর্ণ নোটিশ (\lat{Important} \lat{Notice}):} এই কন্টেন্টটি \lat{AI} (কৃত্রিম বুদ্ধিমত্তা) দিয়ে প্রস্তুত ও সংকলিত। \lat{AI} কঠোর সূত্র-মান নীতি মেনে শুধু নির্ভরযোগ্য উৎস থেকে তথ্য সংগ্রহ করেছে — কেবল সহীহ/হাসান হাদিস ও সহীহ সনদের আসার রাখা হয়েছে, দুর্বল সনদ বাদ দেওয়া হয়েছে। তবে এটি কোনো আলেম বা মুহাদ্দিস কর্তৃক যাচাইকৃত নয়।
\end{warnbox}


\medskip\hrule\medskip

\subsection*{১. সূত্র কার্ড}

\begin{table}[h]\centering\small
\begin{tabular}{ll}
বিষয় & বিবরণ \\
\hline
\textbf{বর্ণনাকারী} & তারিক ইবনে শিহাব (রহ.) থেকে \\
\textbf{গ্রন্থ} & আল-মুস্তাদরাক আলা-স-সহীহাইন (হাকিম), হাদিস ২০৭ ও ২০৮ \\
\textbf{অধ্যায়} & মুস্তাদরাক — সাহাবাদের মর্যাদা \\
\textbf{সনদের মান} & \textbf{সহীহ} (ইমাম হাকিম ও যাহাবির মতে বুখারি-মুসলিমের শর্তানুযায়ী) \\
\textbf{মূল বিষয়} & ইসলামই সম্মানের উৎস — পোশাক-পদ নয়; বিনয়ের সাথে বিজয় \\
\hline
\end{tabular}
\end{table}

\medskip\hrule\medskip

\subsection*{২. পটভূমি (কে, কখন, কোথায়, কেন)}

হিজরি ১৬ সনে মুসলিম সেনারা \textbf{শাম (সিরিয়া-ফিলিস্তিন)} অঞ্চলে \lat{victory} (বিজয়) লাভ করছিল। সেনাপতি \textbf{আবু উবায়দাহ ইবনুল জাররাহ (রাঃ)}-এর নেতৃত্বে মুসলিমরা সেখানে পৌঁছেছে। এলাকার খ্রিস্টান নেতারা শান্তিচুক্তির শর্ত হিসেবে চেয়েছিল — খলিফা নিজে এসে চুক্তি সম্পন্ন করুন।

\textbf{উমার ইবনুল খাত্তাব (রাঃ)} মদীনা থেকে শামের দিকে \lat{journey} (সফর) শুরু করলেন। পথে একটি নদী/নালা (মাখাযা) পেরোতে হলো। উমার উট থেকে নামলেন, নিজের \textbf{জুতা/মোজা (খুফ)} খুলে কাঁধে রাখলেন, উটের লাগাম ধরলেন এবং নিজেই উটকে নিয়ে নদী পাড়ি দিলেন।

বিজয়ী সেনারা — সুসজ্জিত ও সম্মানিত — খলিফাকে দেখলেন সম্পূর্ণ সাধারণ অবস্থায়। সৈন্য ও এলাকার বড় বড় নেতারা (বাতারিকা) তাঁকে সম্বর্ধনার জন্য অপেক্ষা করছিলেন, আর খলিফা নিজেই উট হাঁটাচ্ছেন! এ \lat{scene} (দৃশ্য) দেখে সেনাপতি আবু উবায়দাহ (রাঃ) মন্তব্য করলেন।

\medskip\hrule\medskip

\subsection*{৩. পূর্ণ ঘটনার বর্ণনা}

তারিক ইবনে শিহাব (রহ.) বলেন:

\begin{quote}
\textbf{\lat{“}উমার ইবনুল খাত্তাব (রাঃ) শামের দিকে বের হলেন, আর আমাদের সাথে ছিলেন আবু উবায়দাহ ইবনুল জাররাহ (রাঃ)। তারা একটি নদী/নালা (মাখাযা)-র কাছে পৌঁছালেন — উমার তখন উটের উপর।} \\
\textbf{তিনি উট থেকে নেমে তাঁর জুতা/মোজা (খুফ) খুলে কাঁধে রাখলেন, উটের লাগাম ধরে উটসহ নদী পাড়ি দিলেন।} \\
\textbf{আবু উবায়দাহ (রাঃ) বললেন: "হে আমীরুল মু'মিনীন! আপনি কি এমন করছেন? আপনি নিজের জুতা খুলে কাঁধে রেখেছেন, উটের লাগাম ধরেছেন এবং উটসহ নদী পার হচ্ছেন?! আমি তো চাইতাম না যে শহরের মানুষ আপনাকে এ অবস্থায় দেখুক।"} \\
\textbf{উমার (রাঃ) বললেন: "আহ! হে আবু উবায়দাহ! এ কথা যদি আপনি ছাড়া অন্য কেউ বলত! আমি তো এটা করেছি উম্মতে মুহাম্মদের (সাল্লাল্লাহু আলাইহি ওয়া সাল্লাম) জন্য দৃষ্টান্ত (নাকাল) হিসেবে। আমরা ছিলাম সবচেয়ে লাঞ্ছিত জাতি — অতঃপর আল্লাহ আমাদের ইসলাম দিয়ে সম্মানিত করেছেন। এখন আমরা যদি ইসলাম ছাড়া অন্য কিছুর মাধ্যমে সম্মান চাই, তবে আল্লাহ আমাদের লাঞ্ছিত করবেন।"\lat{”}}
\end{quote}

\subsubsection*{দ্বিতীয় বর্ণনা (হাকিম ২০৮)}

\begin{quote}
\textbf{\lat{“}উমার যখন শামে পৌঁছলেন, সৈন্যরা তাঁর সাথে দেখা করল — তখন তাঁর পরনে ছিল একটি ইযার (তহবন্দ), খুফ (জুতা) ও পাগড়ি; তিনি নিজের উটের মাথার লাগাম ধরে নদী পার হচ্ছিলেন।} \\
\textbf{কেউ বলল: "হে আমীরুল মু'মিনীন! সৈন্যরা ও শামের বড় বড় নেতারা আপনাকে এমন অবস্থায় সম্বর্ধনা করছে?!"} \\
\textbf{উমার (রাঃ) বললেন: "আমরা এমন এক জাতি, যাদের আল্লাহ ইসলাম দিয়ে সম্মানিত করেছেন। আমরা তাই কখনো ইসলাম ছাড়া অন্য কিছুর মাধ্যমে সম্মান খুঁজব না।"\lat{”}}
\end{quote}

\medskip\hrule\medskip

\subsection*{৪. ধাপে ধাপে বিশ্লেষণ}

\textbf{ধাপ ১ — বিজয়ের মুহূর্তে বিনয়:} বিজয়ী খলিফা বিজয়ের উল্লাসে অহংকারী হননি; বরং নিজেই উট হাঁটিয়ে নদী পার হয়েছেন — এমন \lat{humility} (বিনয়) বিরল।

\textbf{ধাপ ২ — সেনাপতির মন্তব্য:} আবু উবায়দাহ (রাঃ) — যিনি উমারের স্নেহভাজন ও বড় সাহাবি — খলিফার এই \lat{appearance} (অবস্থা) দেখে বললেন: \lat{“}আমি চাইতাম না শহরের মানুষ আপনাকে এভাবে দেখুক।\lat{”}

\textbf{ধাপ ৩ — উমারের চমৎকার জবাব:} উমার বললেন — আমি এটা করেছি উম্মতের জন্য \textbf{দৃষ্টান্ত} হিসেবে। অর্থাৎ তিনি \lat{intentionally} (ইচ্ছাকৃতভাবে) এ কাজ করেছেন, যাতে উম্মত শেখে — সম্মান পোশাক-পদে নয়, ইসলামে।

\textbf{ধাপ ৪ — ঐতিহাসিক সত্য:} উমার বললেন — আমরা ছিলাম সবচেয়ে লাঞ্ছিত জাতি; ইসলাম আমাদের \lat{honour} (সম্মান) দিয়েছে। যারা ইসলাম ছাড়া অন্য সম্মান চায়, আল্লাহ তাদের অপমানিত করেন।

\textbf{ধাপ ৫ — চিরন্তন শিক্ষা:} বিজয়, ক্ষমতা ও সম্মানের মুহূর্তে অহংকার এড়িয়ে বিনয় ধরা — এটাই মুমিনের \lat{character} (চরিত্র)। ইসলামই আমাদের গর্বের একমাত্র উৎস।

\medskip\hrule\medskip

\subsection*{৫. শব্দের ব্যাখ্যা}

\begin{table}[h]\centering\small
\begin{tabular}{ll}
শব্দ & অর্থ \\
\hline
\textbf{নাকাল (\arab{نكال})} & দৃষ্টান্ত/শিক্ষা — যার মাধ্যমে অন্যদের শেখানো হয় \\
\textbf{মাখাযা (\arab{مخاضة})} & নদীর অগভীর স্থান যেখানে উট/ঘোড়া দিয়ে পার হওয়া যায় \\
\textbf{খুফ (\arab{خف})} & চামড়ার মোজা/জুতা \\
\textbf{বাতারিকা (\arab{بطارقة})} & শামের বড় বড় নেতা/সেনাপতি (রোমান পদবি) \\
\textbf{আযিল্লাহ (\arab{أذلة})} & লাঞ্ছিত/অপমানিত \\
\textbf{ইজ্জত (\arab{عزة})} & সম্মান/মর্যাদা \\
\hline
\end{tabular}
\end{table}

\textbf{উমারের বাণীর মূল:} ইসলামের আগে আরবরা ছিল সাম্রাজ্যের নিগৃহীত বেদুইন; ইসলাম তাদের এমন সম্মান দিল যে দুই পরাশক্তির বিরুদ্ধে তারা বিজয়ী হলো। এই সম্মান আল্লাহর দান — তাই তা অন্য কোথাও খোঁজা কুফরি ও অহংকার।

\medskip\hrule\medskip

\subsection*{৬. সংশ্লিষ্ট হাদিস ও আয়াত}

\textbf{আয়াত (সূরা আলে ইমরান ৩:১৩৯):}
\begin{quote}
\textbf{\lat{“}তোমরা হতোদ্যম হয়ো না, দুঃখ করো না; যদি মুমিন হও, তবে তোমরাই বিজয়ী।\lat{”}}
\end{quote}

\textbf{আয়াত (সূরা আল-মুনাফিকুন ৬৩:৮):}
\begin{quote}
\textbf{\lat{“}সম্মান তো আল্লাহর, তাঁর রাসূলের ও মুমিনদের।\lat{”}}
\end{quote}

\textbf{আয়াত (সূরা ফাতির ৩৫:১০):}
\begin{quote}
\textbf{\lat{“}যে সম্মান (ইজ্জত) চায়, সে জেনে রাখুক — সম্মান সম্পূর্ণ আল্লাহরই।\lat{”}}
\end{quote}

\textbf{হাদিস (তিরমিজি ২০২৯):} নবীজি বলেছেন: \textbf{\lat{“}যে আল্লাহর জন্য বিনয় করে, আল্লাহ তাকে উন্নত করেন।\lat{”}}

\textbf{হাদিস (মুসলিম ৯১):} কিবর — সত্যকে তুচ্ছ করা ও মানুষকে হেয় করা। উমার বিজয়ে গর্ব না করে বিনয়ী থেকেছেন — কিবর নয়, বরং আল্লাহর সম্মানের মালিকানা স্মরণ করেছেন।

\medskip\hrule\medskip

\subsection*{৭. গভীর শিক্ষা}

\textbf{১. সম্মানের একমাত্র উৎস — আল্লাহ ও ইসলাম।}
উমারের বাণী চিরকালীন: \lat{“}আমরা এমন জাতি যাদের আল্লাহ ইসলাম দিয়ে সম্মানিত করেছেন; ইসলাম ছাড়া অন্য সম্মান চাইব না।\lat{”} — মুমিনের গর্বের \lat{foundation} (ভিত্তি) হওয়া উচিত ঈমান, ধন-পদ নয়।

\textbf{২. বিজয়ের মুহূর্তে অহংকার — সবচেয়ে বড় পরীক্ষা।}
\lat{Success} (সফলতা) ও বিজয় অহংকারের সবচেয়ে শক্তিশালী উৎস। উমার বিজয়ী হয়েও সাধারণ পোশাকে নিজে উট হাঁটিয়েছেন — অহংকার ভাঙার সর্বোত্তম \lat{example} (উদাহরণ)।

\textbf{৩. নেতার জন্য দৃষ্টান্ত হওয়া।}
উমার বলেছেন: \lat{“}আমি এটা করেছি উম্মতের জন্য দৃষ্টান্ত হিসেবে।\lat{”} — নেতা-শাসকদের উচিত এমন \lat{life} (জীবন) যাপন যাতে উম্মত তাদের থেকে বিনয় শেখে।

\textbf{৪. দুনিয়ার সম্মানের মোহ ত্যাগ।}
শামের নেতারা খলিফাকে পোশাক-আভিজাত্যে দেখতে চেয়েছিল; উমার দেখালেন — প্রকৃত \lat{dignity} (মর্যাদা) মানুষের চোখে নয়, আল্লাহর কাছে।

\textbf{৫. ইসলামের সম্মান রক্ষা।}
যারা ইসলাম ছাড়া অন্য কিছু (বংশ, \lat{wealth} (ধন), ক্ষমতা) দিয়ে সম্মান চায়, আল্লাহ তাদের লাঞ্ছিত করেন। মুমিনের উচিত ইসলামকেই গর্বের ভিত্তি রাখা।

\medskip\hrule\medskip

\subsection*{৮. জীবনে প্রয়োগের পদ্ধতি}

\begin{enumerate}
\item \textbf{সফলতার মুহূর্তে বিনয়:} পরীক্ষায় ভালো \lat{result} (ফল), চাকরি, ব্যবসায় সফলতা — এই মুহূর্তে উমারের মতো বিনয় ধারণ করুন, অহংকার নয়।
\item \textbf{পোশাক-পদে গর্ব নয়:} দামি পোশাক, গাড়ি বা \lat{status} (পদমর্যাদা) নিয়ে গর্ব করবেন না; মনে রাখুন — সম্মান ইসলাম ও তাকওয়ায়।
\item \textbf{উম্মতের দৃষ্টান্ত হোন:} নিজের \lat{behaviour} (আচরণ) এমন রাখুন যাতে ছোটরা আপনার থেকে বিনয় শেখে।
\item \textbf{নিজের সম্মানের উৎস চেনা:} \lat{Self-reflection} (আত্ম-পর্যালোচনা) করুন — আমি নিজেকে যেসব বিষয়ে বড় মনে করি, তা কি আল্লাহর দান নাকি আমার নিজের কৃতিত্ব?
\item \textbf{অন্যের সম্মান রক্ষা:} নিজের \lat{humility} বজায় রেখে অন্যদের সম্মান করুন।
\end{enumerate}
\medskip\hrule\medskip

\subsection*{৯. রেফারেন্স}

\begin{itemize}
\item আল-মুস্তাদরাক আলা-স-সহীহাইন (হাকিম), হাদিস ২০৭ ও ২০৮
\item যাহাবির মূল্যায়ন: হাদিস ২০৭ — বুখারি-মুসলিমের শর্তানুযায়ী সহীহ
\item সূরা আলে ইমরান ৩:১৩৯; সূরা আল-মুনাফিকুন ৬৩:৮; সূরা ফাতির ৩৫:১০
\item সহীহ মুসলিম ৯১; সুনানে তিরমিজি ২০২৯
\item ইসলামের ইতিহাস: শাম বিজয় ও উমারের সফর (তারিখুত তাবারী)
\end{itemize}

\end{document}
